\subsection{Análisis de Sensibilidad Bidimensional y Estabilidad Presupuestaria}\label{anuxe1lisis-de-sensibilidad-bidimensional-y-estabilidad-presupuestaria}

\subsection{Sensibilidad por Bandas Salariales del Formulario E-26}\label{sensibilidad-por-bandas-salariales-del-formulario-e-26}

Las dos variables de mayor impacto en la estructura presupuestaria son la tarifa horaria de los perfiles profesionales (Formulario E-26) y la arquitectura tecnológica de custodia de claves:

\begin{enumerate}
\def\labelenumi{\arabic{enumi}.}
\tightlist
\item
  \textbf{Sensibilidad por Bandas E-26:} Variando las tarifas de los roles entre el límite inferior y superior del Formulario E-26, el presupuesto fluctúa entre \textbf{6.812,0 UF} (escenario de costo mínimo) y \textbf{10.150,0 UF} (escenario de tarifa máxima de mercado).
\item
  \textbf{Sensibilidad Criptográfica (Hardware dedicado vs.~Claves individuales):} Si la arquitectura adopta un clúster exclusivo de \textbf{Managed HSM Standard B1} (\(\text{USD } 3,20\text{/h}\) de lista) en lugar de claves protegidas por HSM en Key Vault Premium (\(\text{USD } 1\text{/clave/mes}\)), el costo de gestión de claves se incrementa de \(787,16\text{ UF}\) a \(2.680,6\text{ UF}\), situando el presupuesto total en \textbf{10.385,6 UF} (\(VAN = 8.125,4\text{ UF}\)).
\end{enumerate}

Sobre esa base, la variación conjunta del retainer del DPO (\(\pm 20\%\)) y de la plataforma GRC (\(\pm 25\%\)) mueve el presupuesto dentro de una banda de \(\pm 5{,}57\%\) ---entre 8.018,96 y 8.965,36 UF---, lo que confirma que ninguna de las dos variables de mayor incertidumbre compromete la viabilidad económica del contrato.

\subsection{Racionalidad Financiera: Regla de Gordon-Loeb, RoSI y Umbral de Indiferencia}\label{racionalidad-financiera-regla-de-gordon-loeb-rosi-y-umbral-de-indiferencia}

Bajo la Ley 21.719 (Art. 46), el régimen para infracciones gravísimas contempla multas de hasta \textbf{20.000 UTM} (\(35.000\text{ UF}\) o \(\$1.400\text{ millones de CLP}\)), a las cuales se suma la potestad punitiva de la Ley 21.663 (hasta 10.000 UTM) y los costos de remediación forense DFIR, totalizando una exposición contingente agregada de \textbf{55.000 UF} (\(\$2.200\text{ millones de CLP}\)).

De acuerdo con el modelo económico de \textbf{Gordon y Loeb (2002)}, la inversión óptima en ciberseguridad y protección de datos se acota a un techo del \(37\%\) de la pérdida esperada: {\small\[\text{Presupuesto Óptimo} \le 0,37 \times \text{Pérdida Esperada} = 0,37 \times 55.000\text{ UF} = \mathbf{20.350\text{ UF}}\]} El costo total del programa AudIT (\(8.492,16\text{ UF}\)) representa solo el \textbf{\(15,4\%\) de la exposición patrimonial}, situándose holgadamente bajo la cota de sobreinversión.

Evaluando el Retorno sobre la Inversión en Seguridad (\(\text{RoSI}\)): {\small\[\text{RoSI} = \frac{(\text{Exposición Punitiva} \times \text{Eficacia Mitigación}) - \text{TCO}}{\text{TCO}} \times 100\% = \frac{(55.000\text{ UF} \times 0,85) - 8.492,16\text{ UF}}{8.492,16\text{ UF}} \times 100\% = \mathbf{450,51\%}\]} El \textbf{umbral de probabilidad de indiferencia} es: {\small\[p^* = \frac{\text{TCO}}{\text{Exposición Punitiva}} = \frac{8.492,16\text{ UF}}{55.000,0\text{ UF}} = \mathbf{15,44\%\text{ en 56 meses}} \implies \mathbf{3,05\%\text{ anual}}\]} Basta con que la probabilidad anual de sufrir un incidente sancionable supere el \textbf{\(3,05\%\)} para que el programa de cumplimiento genere un beneficio económico neto directo para Transportes Curimón S.A., blindando el flujo de caja del consorcio adjudicatario.
